\chapter{反应为什么有快有慢}

\begin{kaichang}
你也许在生日聚会或者庙会上玩过荧光棒：一根软软的塑料管，摸上去什么动静都没有。用力“啪”地一掰，摇一摇，它就亮了——幽幽的绿光或蓝光，不用电池，不插电，摸上去也不烫。

现在做两个小小的折腾。把它丢进一杯热水里，它会猛地变亮，亮得精神十足；可是不到半小时，就暗下去了。再拿一根新的放进冰箱冷藏室，它会亮得蔫蔫的，却能撑上一整夜。同一根荧光棒，亮度仿佛被一只看不见的旋钮拧来拧去。

先猜一猜：热水到底“拧”动了什么？是让发光反应变多了，还是变快了？这有区别吗？这一章结束时，你会亲手拧动这只旋钮。
\end{kaichang}

\section{从一声爆炸到千年风化}

先收集一堆你早就见过的事实。鞭炮炸开，是眨眼之间的事；泡腾片丢进水里，几十秒就安静下来；一瓶鲜牛奶在冰箱里能放几天，在夏天的灶台上半天就变味；铁钉生锈要几个星期；一个塑料袋埋在土里要几百年才烂得差不多；山里的岩石被风雨剥蚀，动辄几千年；溶洞里的钟乳石长高一厘米，可能要上百年。

\begin{table}[htbp]
\centering
\begin{tabular}{ll}
\toprule
\textbf{过程} & \textbf{大致时间尺度}\\
\midrule
酸碱相遇、炸药爆炸 & 一瞬间（微秒到毫秒）\\
泡腾片冒泡 & 几十秒\\
荧光棒发光 & 几小时\\
牛奶变质、苹果切面发黄 & 几小时到几天\\
铁钉生锈 & 几周到几年\\
塑料袋在土壤中降解 & 几百年\\
岩石风化、钟乳石生长 & 上千年到上万年\\
\bottomrule
\end{tabular}
\end{table}

同样是“化学反应”，时间跨度竟从不到一秒到上万年——差了十几个数量级。更值得注意的是：\textbf{同一个反应，也能忽快忽慢}。铁在干燥空气里慢慢长斑，在潮湿的盐雾中锈得飞快，在纯氧里点燃铁丝却能剧烈燃烧、火星四溅。再看厨房：鸡蛋扔进沸水，几分钟就全熟；放进约 $63\,^{\circ}\mathrm{C}$ 的温水里泡半小时，切开还是溏心——同样是让蛋白质变性凝固，温度差不到 $40\,^{\circ}\mathrm{C}$，速率差出几十倍。第 27 章讲过反应放热还是吸热，第 28 章讲过自由能决定反应朝哪个方向走——可“方向”和“放多少热”都回答不了本章的问题：\textbf{它到底跑多快}。方向归热力学管，速度归另一套规则管，这套规则叫\textbf{动力学}。这正是本章要拆的机器。

\section{反应先要从一次碰撞开始}

把镜头推进到粒子层。第 26 章说过，化学方程式只是一本“原子重新组合”的账本：旧分子拆开，原子换搭档，新分子出场。可是，两个分子凭什么换搭档？第一步几乎总是——\textbf{它们得先撞上}。原子和分子没有手，不会互相寻找，它们只能靠永不停歇的热运动瞎碰（第 9 章）。没人撞上，一切都免谈。

那么，撞上就能反应吗？算一笔账你就知道不能。常温常压下，一升空气里大约有 $2.5\times 10^{22}$ 个分子，每个分子以每秒几百米的速度飞奔，平均每秒要和别的分子撞上大约 $10^{9}$ 到 $10^{10}$ 次。假如每次碰撞都有效，一升气体里的反应物会在远小于百万分之一秒的时间里全部耗尽——可现实是，氢气和氧气装在一个瓶子里，放上几年也几乎不动。结论只能是：\textbf{绝大多数碰撞都是“无效碰撞”}，两个分子碰一下，又各自弹开，什么也没发生。反应的快慢，本质上是“天文数字的碰撞次数”与“小得可怜的合格率”之间的一场拔河。

什么样的碰撞才算\textbf{有效碰撞}？两个条件，缺一不可。

第一，\textbf{撞得要够狠}。旧键要松开、原子要换位置，先要跨过一道能量门槛。慢悠悠蹭过来的两个分子，挤一挤就弹开了；只有带着足够大动能迎面撞上的，才可能把旧结构“撞松”。这道门槛的最低高度，就是大名鼎鼎的\textbf{活化能}（记作 $E_{\mathrm{a}}$）。

第二，\textbf{撞的要害部位要对}。分子不是圆球，有形状、有方向。比如一个氢分子要撞上碘分子让 \chem{H—I} 键搭上关系，撞在碘分子的“侧面”才有戏，撞在别的角度就只是擦身而过。取向不对的碰撞，能量再大也白搭。

把这两条合起来，反应速率的微观画像就出来了：\textbf{速率 = 碰撞次数 × 有效碰撞所占的比例}。想让反应变快，要么让它们撞得更勤，要么让每次碰撞更有可能“合格”。后面所有的调速旋钮，拧的都是这两项。

\section{一座必须翻过的山}

把“能量门槛”画出来，就是化学里最重要的一张图。横轴是反应进程（从反应物到生成物），纵轴是能量：

\begin{center}
\begin{tikzpicture}[scale=0.9]
\draw[->] (0,0) -- (9,0) node[right]{\footnotesize 反应进程};
\draw[->] (0,0) -- (0,4.8) node[above]{\footnotesize 能量};
\draw[thick] plot[smooth] coordinates {(0.5,2.8) (2.2,2.9) (4,4.2) (5.8,2.0) (7.5,1.2)};
\draw[dashed] (0.5,2.8) -- (2.0,2.8);
\draw[dashed] (7.5,1.2) -- (8.3,1.2);
\draw[<->] (4.7,2.8) -- (4.7,4.15) node[midway,right]{\footnotesize 活化能 $E_{\mathrm{a}}$};
\node at (1.2,3.25){\footnotesize 反应物};
\node at (7.9,0.75){\footnotesize 生成物};
\node at (4,4.6){\footnotesize 过渡态};
\end{tikzpicture}

\footnotesize 能量示意图（人为表示法）：曲线不是分子真的爬的山路，只是“反应进行中体系能量怎样变化”的记账方式。
\end{center}

反应物在山这边，生成物在山那边。这张图里，生成物的位置比反应物低——所以这是个放热反应（第 27 章）。但请注意：\textbf{终点再低，也要先翻过山顶}。山顶上那个半旧不新、键断了一半又接上了一半的短命结构，叫\textbf{过渡态}。它有点像两个人传球时球恰好悬在两手之间的那个瞬间：既不是甲拿着，也不是乙拿着，却是整个交接中最关键、也最难“抓拍”的一刹那。过渡态只存在约 $10^{-13}$ 秒量级，比一个分子振动一次的时间长不了多少，却决定了整个反应的快慢：山越矮，一次碰撞就越容易“合格”，反应越快；山越高，分子们撞得鼻青脸肿也过不去，反应再“该发生”也只能干等。

这张图顺手澄清了前两章留下的一个悬念。第 28 章说，$\Delta G$ 为负的反应“可以自发发生”。很多人想当然地补上一句：那就马上发生呗。错。汽油和氧气混合，$\Delta G$ 是深深的负值，可它们在油箱里和平共处，直到一颗火花出现——火花给了局部一小撮分子翻越活化能的能量，它们反应放热，加热周围更多分子翻山，链式推进，于是轰的一声。\textbf{自发，只保证山那边更低；快不快，要看山有多高。}方向是热力学的事，速度是动力学的事，两本账分开记。

那温度是怎么帮忙的？关键在于：一群分子的能量并不整齐，而是有大有小、呈统计分布——大多数分子挤在平均值附近，少数“高能分子”拖在尾巴上。活化能就像一根分数线：能量在线以上的分子，碰撞才可能合格。升温一度，平均能量只涨一点点，但\textbf{尾巴上“过线分子”的比例会指数式地放大}。这就是为什么温度对速率的影响如此凶猛。

\begin{qiaoliang}
1889 年，阿伦尼乌斯把这种关系写成了一个式子：
\[k = A\,\mathrm{e}^{-E_{\mathrm{a}}/(RT)}\]
其中 $k$ 是速率常数（衡量“合格碰撞转化成反应”的本领），$T$ 是开尔文温度，$R$ 是气体常数，$E_{\mathrm{a}}$ 是活化能。别怕这个式子，只记住一件事：温度 $T$ 在分母上、又躲在指数里，所以它对 $k$ 的影响是\textbf{指数级}的。代入具体数字感受一下：取一个典型的 $E_{\mathrm{a}} = \ud{50}{kJ\,mol^{-1}}$，从 $300\,\mathrm{K}$（约 $27\,^{\circ}\mathrm{C}$）升到 $310\,\mathrm{K}$，
\[\frac{k_{310}}{k_{300}} = \mathrm{e}^{\frac{E_{\mathrm{a}}}{R}\left(\frac{1}{300}-\frac{1}{310}\right)} = \mathrm{e}^{0.65} \approx 1.9\]
速率几乎翻了一倍。这就是厨房里那条经验法则的由来：\textbf{许多反应温度每升高约 $10\,^{\circ}\mathrm{C}$，速率变为原来的 2 到 3 倍}（活化能越大，倍数越猛）。冰箱冷藏室约 $4\,^{\circ}\mathrm{C}$，室温约 $24\,^{\circ}\mathrm{C}$，相差 $20\,^{\circ}\mathrm{C}$——变质反应在室温下大约快 $2\times 2=4$ 倍。这就是冰箱存在的全部理由：它不改变化学反应的方向，只是把所有反应的“山”集体加高了一截。
\end{qiaoliang}

\section{四个调速旋钮}

有了“速率 = 碰撞次数 × 合格比例”这把钥匙，四个常见因素就能一次讲通。

\textbf{浓度}。一升水里溶质的分子越多，每个分子撞上“对的人”的机会就越大——碰撞次数直接正比于浓度。铁丝在空气里（氧气约占五分之一）只是发红，伸进一瓶纯氧里立刻火星四射，就是同一个道理。医药柜里 $3\%$ 的过氧化氢（双氧水）消毒时慢悠悠冒小泡，而高浓度的过氧化氢能猛烈分解、甚至用作火箭推进剂的成分——浓度这只旋钮，既能调快慢，也能调出完全不同的危险等级。

\textbf{温度}。上一节算过了：升温对碰撞次数的贡献很小（分子速度只增加百分之几），真正的威力在于指数式放大合格碰撞的比例。它是四个旋钮里劲最大的一个。

\textbf{接触面积}。固体和液体反应时，只有表面的粒子能挨打。一根铁钉几天才见锈斑，同样质量的铁粉撒进空气里遇到火星却能瞬间爆燃——表面积差了几千上万倍。把块状物砸碎、把固体搅成悬浊液，都是把更多粒子“暴露到前线”。这条旋钮拧过头会出事：磨得极细的面粉、煤粉悬浮在空气里，表面积大得惊人，一粒火星就能让整个车间变成一个大火药桶。所以面粉厂、煤矿把“严禁烟火”写进铁律，这不是迷信，是动力学。

\textbf{压强}。只对气体反应管用：压缩气体等于把分子挤得更密，单位体积里的碰撞次数随之上升。工厂里合成氨要在高压下进行，高压既是让分子们“挤着见面”（速率），也逼着反应往分子数少的方向走（平衡，第 31 章细讲）。压力锅则是拧了另一个旋钮：锅内压强高，水的沸点从 $100\,^{\circ}\mathrm{C}$ 升到约 $120\,^{\circ}\mathrm{C}$，按“每升 $10\,^{\circ}\mathrm{C}$ 快两三倍”估算，炖肉的反应快了四五倍，于是软烂得更快。

\begin{shiyan}
\textbf{泡腾片计时赛（在家可做）}

材料：泡腾片 4 片、透明杯子 4 个、冷水、温水（手感温热即可）、热水（请家长倒，$60\,^{\circ}\mathrm{C}$ 以内）、计时器、勺子。

步骤：① 三个杯子装等量的冷水、温水、热水，同时各投入一整片泡腾片，记下每杯“气泡基本平息”的时间。② 第四杯装与某一杯同温度的水，把一片泡腾片碾碎后撒入，计时对比。

观察点：哪一杯气泡最猛、结束最早？整片和粉末差多少？注意别只看“反应剧烈”，要记“反应结束的时间”——剧烈是因为快，结束早也是因为快，这是同一回事的两张脸。

安全提示：热水由大人操作，谨防烫伤；\textbf{绝对不要}把泡腾片和水装进密封瓶里摇晃——产生的气体可能让瓶盖像炮弹一样飞出去。
\end{shiyan}

\section{速率方程：只能从实验里读出来的账}

现在轮到符号登场。化学家用\textbf{速率方程}给反应的快慢记账，比如对反应 $\chem{A} + \chem{B} \rightarrow \chem{C}$，速率常常写成
\[v = k \cdot c(\chem{A})^{m} \cdot c(\chem{B})^{n}\]
读作：反应速率 $v$ 与各反应物浓度的若干次方成正比，$k$ 就是上节见过的速率常数。请注意那两个指数 $m$ 和 $n$——它们\textbf{只能从实验测出来，不能照着配平的化学方程式抄}。

为什么不能抄？因为配平方程只记总账：多少个 A 加多少个 B，变成多少个 C（第 26 章）。它完全不记过程。真实反应几乎都不是“一次撞成”，而是分好几步走，像一条流水线。流水线有多快，由\textbf{最慢的那道工序}决定——化学里叫\textbf{决速步}。浓度对速率的影响，取决于谁参与了那道最慢的工序，而不是谁在总账上露面。

有个经典例子能把这个道理钉死。二氧化氮和一氧化碳的反应，账本上写得明明白白：
\[\chem{NO_2} + \chem{CO} \rightarrow \chem{NO} + \chem{CO_2}\]
可实验测出来的速率方程却是 $v = k \cdot c(\chem{NO_2})^{2}$——速率只跟 \chem{NO_2} 的浓度有关，多加一氧化碳，速率几乎纹丝不动！合理的解释是：这反应分两步走，先由两个 \chem{NO_2} 分子慢吞吞地撞出一个 \chem{NO_3}（决速步），\chem{NO_3} 再飞快地把氧塞给 \chem{CO}。一氧化碳只在第二道工序出现，它的浓度自然管不了整条流水线的快慢。反过来，三个及以上粒子“同时撞在一起”的概率小到可以忽略，所以配平系数动辄为 3、4 的反应，几乎必然是多步完成的。账本是账本，流水线是流水线。

那实验上究竟怎么“测”出一个速率？思路出奇地朴素：盯住一个随反应变化的量，看它随时间变得多快。冒气体的反应，可以数单位时间收集到的气泡体积；有颜色变化的反应（比如碘遇淀粉变蓝），可以盯颜色加深的快慢；气体分子数改变的反应，密闭容器里的压强本身就是一只计时器。把浓度换算进去，就得到速率。换几种起始浓度各测一遍，速率方程的指数 $m$、$n$ 就现形了——浓度加倍，速率若变成 2 倍，指数就是 1；若变成 4 倍，指数就是 2；若根本不动，指数就是 0。这套“只改变一个条件、其余全部摁住”的做法，你在本章末的思考题里会亲手设计一遍。

\section{科学家怎样“看见”只存在一瞬间的过渡态}

\begin{zhengju}
过渡态这东西，从提出那天起就带着点“信则有”的味道：它存在的时间大约只有 $10^{-13}$ 秒，没有任何容器装得住它，凭什么说它真的存在？很长一段时间里，化学家手里只有间接证据：速率方程怎么随浓度变化、温度怎么影响速率、把分子里的氢换成更重的同位素速率变了多少……这些线索拼起来，都指向“有一个能量最高、寿命极短的中间结构”，但谁也没亲眼见过。

1980 年代，埃及裔美国化学家泽维尔（Ahmed Zewail）想到了一个简单粗暴的思路：\textbf{给它拍照}。难点在于相机快门必须比 $10^{-13}$ 秒还快。他的办法是用激光：激光脉冲可以短到只有几十飞秒（1 飞秒 $= 10^{-15}$ 秒），比过渡态的寿命还短上千倍。实验设计像用两盏频闪灯拍一段舞蹈：第一束激光（泵浦光）像发令枪，把分子“踢”上反应的山坡；隔上精确控制的几十飞秒，第二束激光（探测光）打过来，照一照此刻分子的“模样”。不断改变两束光之间的时间差，就能得到一串“分子电影”：旧键松了、过渡态出现了、新键成了。1987 年，他的团队用这个方法拍下了氰化碘（\chem{ICN}）分裂成 \chem{I} 和 \chem{CN} 的全过程，过渡态的信号和理论预言的位置、寿命对得上。会不会那信号其实是别的什么中间体？研究团队用不同波长的探测光交叉验证，并和量子力学计算逐点比对，排除了这种替代解释。1999 年，泽维尔独自获得诺贝尔化学奖，飞秒化学从此成为一门学科。这个故事的要点是：连“只存在百万分之一微秒的东西”，科学家也不靠相信，而是靠造出比它更快的快门去验证。
\end{zhengju}

\section{这台模型什么时候会失灵}

碰撞理论简单好用，但它的适用边界要心里有数。它默认粒子像台球一样均匀乱撞，这\textbf{对气体和稀溶液里的简单反应}很管用。可是：有的分子不撞任何东西，自己受热也会分解（单分子反应），能量是它早就储存在内部的；在黏稠的液体里，两个分子撞上后被周围的“溶剂笼子”困住，反复碰又反复弹开，碰撞次数的概念要重新修订；有的反应快得惊人，每一步碰撞都有效，速率的瓶颈变成“粒子游过来见面的速度”，这类反应叫扩散控制，升温对它们的作用小得多。还有一些反应的能量根本不走“碰撞”这条路：照相底片感光、皮肤被太阳晒伤、绿叶捕捉阳光（第 54 章），都是光子把能量直接“快递”进分子内部，撞上没撞上都无所谓，这时要请出光化学的模型。还有更实际的一条：速率方程是在某个温度、浓度范围内测出来的\textbf{经验规律}，把它外推到没测过的区间——比如拿室温数据去算发动机燃烧室里的速率——经常会错得离谱。模型在哪片地上好用、在哪片地上会塌，和使用模型同样重要。

\begin{wuqu}
“放热越多的反应，一定越快。”——错，而且错得很典型。放不放热，看的是起点和终点的高度差；快不快，看的是山顶的高度。这是两座互不相干的“山”。反例俯拾皆是：铁生锈是放热反应，焓变不小，却慢到要用年月计；纸张燃烧放出大量热，可纸在空气里摆上几十年也不着火，因为活化能这座山太高。反过来，有些吸热反应（比如小苏打和柠檬酸）搅在一起立刻冒泡，快得很。判断快慢，问活化能；判断放热，问焓变；判断方向，问自由能。三个问题，三本账，别串。
\end{wuqu}

\section{回到那根荧光棒}

现在来兑现开场的承诺。荧光棒的塑料管里其实装着两拨化学品：一拨溶在管子里，另一拨封在一根细细的玻璃管里，彼此隔开，相安无事。你“啪”地一掰，玻璃管碎裂，两种液体混在一起，反应开始。这个反应释放的能量不像燃烧那样变成热，而是直接交给染料分子，再由染料以光的形式吐出来——所以它亮而不烫，人称“冷光”。

那只看不见的旋钮，现在你有名字了：\textbf{温度}。丢进热水，分子的热运动加剧，更关键的是“能量超过活化能”的分子比例指数式上升，有效碰撞猛增——同一秒内有更多分子翻过山顶，光自然更亮。但管里的反应物总量是固定的，翻山翻得快，存货就烧得快，所以亮得耀眼，也死得干脆。放进冰箱则相反：反应被集体踩了刹车，亮度低了，有限的存货却省着用，能撑一整夜。亮度是速率，续航是总量，你拧温度旋钮时，其实是在“亮得短”和“暗得久”之间做交易——一笔纯粹的动力学交易。下次再掰亮一根荧光棒，你看到的就不再是一根玩具，而是一支握在手里、快慢可调的化学反应。

\section{把快慢做成技术，以及一座能被削矮的山}

人类和反应速率斗了几千年，如今这套知识已经长成一个庞大的产业。食品工业是最熟悉的战场：冷链物流让生鲜在低温中“慢速播放”；巴氏杀菌则走另一条路——把牛奶加热到约 $72\,^{\circ}\mathrm{C}$ 保持十几秒，让细菌体内的蛋白质变性反应快跑到足以杀灭大部分细菌，又不至于把牛奶的风味反应一起烧出来；罐头更狠，$121\,^{\circ}\mathrm{C}$ 高压蒸汽直接“快进到全灭”。充氮包装抽走氧气，是从“浓度”旋钮上动手；厨房里的研磨、切碎、搅拌，全是在拧“接触面积”——火力发电厂把煤磨成比面粉还细的煤粉喷进炉膛，让燃烧在零点几秒内完成，拧的也是它。做酸奶要恒温在约 $42\,^{\circ}\mathrm{C}$，是让乳酸菌的代谢反应跑在最舒服的挡位上，温度低了菌“磨洋工”，温度高了菌被烫死。医生用的许多药物说明书上写“避光、密封、冷藏保存”，翻译过来就是：请帮这个分子把活化能的山加高、把碰撞的机会降到零。

可你有没有觉得哪里不对劲？加热费能源，冷藏也费能源，人类的整条冷链每年烧掉巨量电力，只为给反应踩刹车；工厂为了速率把反应釜烧到几百摄氏度，账单同样吓人。这就引出一个诱人的问题：\textbf{升温只是让分子“更使劲地撞山”——有没有办法直接把山本身削矮？}让反应在常温下也能快跑？有的。而且你身体里此刻就有几万亿台这样的机器在日夜工作，没有它们，你吃下去的午饭要几十年才消化得完。它们叫催化剂——下一章，我们去拆这台机器，看看它怎样在不动总账的前提下，给反应另开一条隧道。

\begin{tiaozhan}
（可跳过，不影响主线）由阿伦尼乌斯式出发，在两个温度 $T_1$、$T_2$ 下取对数相减，可以得到
\[\ln\frac{k_2}{k_1} = \frac{E_{\mathrm{a}}}{R}\left(\frac{1}{T_1}-\frac{1}{T_2}\right)\]
这个式子透露了“每升 $10\,^{\circ}\mathrm{C}$ 翻倍”的秘密：它其实只在 $E_{\mathrm{a}}$ 约 $50\sim\ud{80}{kJ\,mol^{-1}}$、且在室温附近时碰巧成立。活化能越小，温度越拧不动速率——扩散控制的反应（$E_{\mathrm{a}}$ 只有十几 $\mathrm{kJ\,mol^{-1}}$）每升 $10\,^{\circ}\mathrm{C}$ 只快约 1.2 倍；反之，$E_{\mathrm{a}}$ 极高的反应对温度敏感得近乎神经质。还有一个好玩的推论：如果一个反应“升温反而变慢”，那它多半不是一步反应，而是某一步的平衡被温度挪了位置——这又把话题送回第 31 章了。
\end{tiaozhan}

\section*{练一练}

\begin{sikaoti}
\item 用圆圈和箭头画一幅粒子图：画出分子 A—B 与分子 C 的一次“有效碰撞”和一次“无效碰撞”，在图旁标出两次碰撞在取向或能量上的差别。注明：圆圈大小和箭头粗细都是人为表示法。
\item 画两张能量图：一张是“放热但很慢”的反应，一张是“放热且很快”的反应。两张图的哪些部分必须长得不一样？哪些部分可以一样？
\item 泡腾片实验里，把“碾碎 + 温水”与“整片 + 冷水”两组对比，谁先结束？如果“碾碎 + 冷水”与“整片 + 温水”比呢？用两个调速旋钮分别分析，并说说为什么第二组谁赢要靠实验才知道。
\item 一盒牛奶在 $4\,^{\circ}\mathrm{C}$ 的冷藏室里能放 7 天。按“每升 $10\,^{\circ}\mathrm{C}$ 速率翻倍”估算，它在 $24\,^{\circ}\mathrm{C}$ 的室温下大约能放多久？再想想：真实的牛奶变质主要是细菌繁殖引起的，这个估算哪里可能失真？
\item 化学方程式 $\chem{2HI} \rightarrow \chem{H_2} + \chem{I_2}$ 的配平系数是 2。这是否意味着反应一定由两个 \chem{HI} 分子同时相撞完成？结合“决速步”和“多步反应”谈谈：为什么配平方程不能直接翻译成速率方程？
\item 设计一个实验，测定“某金属与稀酸反应的速率怎样随酸的浓度变化”。你需要测什么量、固定哪些条件（温度、金属表面积、用量……）、改变哪一个条件？画出你预想的结果曲线的大致形状。
\end{sikaoti}
